\documentclass[12pt]{article}
\usepackage[utf8]{inputenc}
\usepackage[T1]{fontenc}
\usepackage{hyperref}
\title{Koropó Language Resources for K-12 Curriculum Development}
\author{Compiled by Claude Sonnet 5.5 (claude-sonnet-5-5) for Andr\'as Kornai}
\date{}
\begin{document}
\maketitle

\section{Introduction}
Koropó (ISO 639-3 xxr) is an extinct language of eastern Brazil, once spoken in Minas Gerais and Rio de Janeiro. It is known only from two early 19th-century word lists made by German travelers: about 127 words (Eschwege, 1815) and about 55 words (Schott, 1818). Its classification is disputed. Silva Neto (2007) treated it as Purian. Nikulin (2020) and Ramirez, Vegini and França (2015) place it in Macro-Jê, as Maxakalían, and explain the similarities with Puri as loanwords. No native speakers, no grammar and no teaching materials were found.

\section{Primary Educational Resources}
No primers, textbooks or school programmes for Koropó turned up in the searches. A curriculum would have to be built from the historical word lists and the scholarly analyses listed below.

\section{Linguistic and Reference Materials}
\begin{itemize}
\item Wikipedia, Koropó language (overview of the sources and the classification debate): \url{https://en.wikipedia.org/wiki/Korop%C3%B3_language}
\item Wikipedia, Purian languages (the related family and the disputed link): \url{https://en.wikipedia.org/wiki/Purian_languages}
\item Wikipedia, Puri language: \url{https://en.wikipedia.org/wiki/Puri_language}
\item Silva Neto (2007), University of Brasília dissertation giving a normalized revision of the Puri, Coroado and Koropó data: \url{https://repositorio.unb.br/bitstream/10482/2328/1/2007_AmbrosioPereiradaSilvaNeto.pdf}
\item Ramirez, Vegini and França (2015), \emph{LIAMES}, on the classification of Koropó: \url{https://periodicos.sbu.unicamp.br/ojs/index.php/liames/article/view/8642302}
\end{itemize}

\section{Religious and Community Resources}
No religious texts or community organisations for Koropó were found in the search results.

\section{Digital and Audio Resources}
No audio, video or software resources were found. Because the language is known only from written word lists, no recordings of it can exist.

\section{Curriculum Development Context}
Very little is documented, so any K-12 programme would be a heritage-awareness programme. It could cover the history of the people, the nature of the sources and the classification debate, and it could use the attested words as a small symbolic vocabulary. Anything beyond that would be reconstruction, and the sources do not support it. Descendant communities in Minas Gerais and Rio de Janeiro should be consulted. No such community was identified in the searches.
\end{document}